\section{Estructura del historial de mensajes de Codex}

\subsection{Cuatro mensajes preinyectados}

Antes de que el usuario emita su solicitud real, Codex ya ha implantado de antemano al menos cuatro mensajes:

\begin{importantbox}{Estructura de mensajes preinyectados de Codex}
\begin{enumerate}
    \item \textbf{Mensaje System}: instrucción básica del sistema (System Prompt)
    \item \textbf{Mensaje Developer (primero)}: inyección de permisos, define los límites de las facultades operativas de Codex
    \item \textbf{Mensaje User (inyectado)}: contiene el contenido de agents.md, el resumen de Skills y el contexto del entorno
    \item \textbf{Mensaje Developer (segundo)}: define el modo de colaboración, describe cómo debe interactuar Codex con el usuario
\end{enumerate}
\end{importantbox}

\subsection{Contenido detallado de cada mensaje}

\subsubsection{Primero: System Instruction}

El mensaje System es la instrucción de más alto nivel, define las normas básicas de conducta y los límites de las capacidades de Codex. Es un System Prompt estándar.

\subsubsection{Segundo: inyección de permisos del Developer}

El mensaje Developer se usa para inyectar la configuración relacionada con los permisos, define los límites de las facultades de Codex en operaciones de archivos, ejecución de comandos y otros aspectos.

\subsubsection{Tercero: inyección del lado User}

\begin{knowledgebox}{El contenido rico del mensaje User}
El tercer mensaje, aunque está marcado con el rol User, en realidad es inyectado automáticamente por la herramienta Codex e incluye tres partes principales:
\begin{itemize}
    \item \textbf{agents.md}: archivo de configuración de Agent a nivel de sistema o de usuario, que se carga desde el directorio \texttt{.codex} o desde la raíz del proyecto de código
    \item \textbf{Resumen de Skills}: enumera todas las Skills disponibles junto con su breve descripción (un índice de carga progresiva)
    \item \textbf{Contexto del entorno}: el directorio de trabajo actual, información del sistema operativo, la fecha actual, la zona horaria y demás información de tiempo de ejecución
\end{itemize}
\end{knowledgebox}

Sobre el registro de las Skills, su estructura incluye:
\begin{itemize}
    \item cuáles Skills están available
    \item la breve descripción de cada Skill
    \item las dos Skills integradas predeterminadas oficiales: Skill Creator y Skill Installer
    \item las instrucciones sobre cómo usar cada Skill
\end{itemize}

\subsubsection{Cuarto: modo de colaboración del Developer}

El cuarto mensaje Developer define cómo debe colaborar Codex con el usuario: estilo de interacción, forma de responder, gestión de Turn, etc. También mantiene una marca/contexto de Turn que se usa para dar seguimiento a los turnos de la conversación.

\subsection{Resumen de la sección}

Antes de que comience la interacción con el usuario, Codex ya implantó 4 mensajes que cubren la instrucción del sistema, los permisos, la configuración de Agent más Skills más entorno, y el modo de colaboración, con lo que construye un contexto completo de Agent Runtime.
